\section*{NeurIPS Paper Checklist}

\begin{enumerate}

\item {\bf Claims}
    \item[] Question: Do the main claims made in the abstract and introduction accurately reflect the paper's contributions and scope?
    \item[] Answer: \answerYes{}
    \item[] Justification: The abstract and introduction state that \coolname{} is a unified RGB-D benchmark with five perception tasks, three-phase capture, and ESD-stratified difficulty splits. Sections~\ref{sec:dataset}--\ref{sec:experiments} directly support each claim with dataset statistics, metric definitions, and experimental results. Scope limitations (indoor-only, D435 depth range, human-driven capture trajectories) are stated in Section~\ref{sec:dataset} and the conclusion.

\item {\bf Limitations}
    \item[] Question: Does the paper discuss the limitations of the work performed by the authors?
    \item[] Answer: \answerYes{}
    \item[] Justification: The conclusion contains a dedicated \emph{Limitations} paragraph discussing T265 odometry drift, D435 depth range constraints, indoor-only capture, and annotator diversity. The dataset card (Croissant) repeats these.

\item {\bf Theory assumptions and proofs}
    \item[] Question: For each theoretical result, does the paper provide the full set of assumptions and a complete (and correct) proof?
    \item[] Answer: \answerNA{}
    \item[] Justification: The paper presents an empirical benchmark and dataset; it contains no theoretical results or proofs.

    \item {\bf Experimental result reproducibility}
    \item[] Question: Does the paper fully disclose all the information needed to reproduce the main experimental results of the paper to the extent that it affects the main claims and/or conclusions of the paper (regardless of whether the code and data are provided or not)?
    \item[] Answer: \answerYes{}
    \item[] Justification: The evaluation toolkit is pip-installable (\texttt{rpx-benchmark}), open-source (MIT), and includes a continuous-integration test suite verifying reproducibility of all metric computations. Sensor specifications, annotation protocol details, and evaluation setup are described in Sections~\ref{sec:dataset} and~\ref{sec:tasks}. All model checkpoints are public.

\item {\bf Open access to data and code}
    \item[] Question: Does the paper provide open access to the data and code, with sufficient instructions to faithfully reproduce the main experimental results, as described in supplemental material?
    \item[] Answer: \answerYes{}
    \item[] Justification: The dataset is hosted on HuggingFace Datasets under CC BY 4.0 with a 5-scene reviewer sample available without registration. The evaluation toolkit source code is publicly available under the MIT licence with pip-installable dependencies and a documented CLI.

\item {\bf Experimental setting/details}
    \item[] Question: Does the paper specify all the training and test details (e.g., data splits, hyperparameters, how they were chosen, type of optimizer) necessary to understand the results?
    \item[] Answer: \answerYes{}
    \item[] Justification: Section~\ref{sec:dataset} describes the train/val/test splits (70/15/15 scenes, no overlap), Section~\ref{sec:features} defines ESD computation and thresholds, and Section~\ref{sec:experiments} reports GPU type, inference settings, and per-model checkpoint identifiers. No training or fine-tuning is performed; all models are evaluated zero-shot.

\item {\bf Experiment statistical significance}
    \item[] Question: Does the paper report error bars suitably and correctly defined or other appropriate information about the statistical significance of the experiments?
    \item[] Answer: \answerYes{}
    \item[] Justification: Tables report standard deviation across scenes alongside mean values. The ESD sensitivity analysis (Appendix~A) uses 1{,}000 Monte Carlo weight perturbations to assess split stability.

\item {\bf Experiments compute resources}
    \item[] Question: For each experiment, does the paper provide sufficient information on the computer resources (type of compute workers, memory, time of execution) needed to reproduce the experiments?
    \item[] Answer: \answerYes{}
    \item[] Justification: Section~\ref{sec:experiments} states GPU type, VRAM, and per-model inference time. The toolkit automatically reports FLOPs and median per-sample latency in every run.

\item {\bf Code of ethics}
    \item[] Question: Does the research conducted in the paper conform, in every respect, with the NeurIPS Code of Ethics \url{https://neurips.cc/public/EthicsGuidelines}?
    \item[] Answer: \answerYes{}
    \item[] Justification: The dataset contains indoor daily objects; no faces, biometric identifiers, or personally identifiable information are captured. All capture was performed by lab personnel who provided informed consent for data release. The work does not raise dual-use concerns.

\item {\bf Broader impacts}
    \item[] Question: Does the paper discuss both potential positive societal impacts and negative societal impacts of the work performed?
    \item[] Answer: \answerYes{}
    \item[] Justification: The conclusion includes a \emph{Broader Impact} paragraph discussing both the positive impact (enabling safer robot deployment through rigorous perception evaluation) and the limited risk profile (indoor daily-object scenes with no misuse pathway).

\item {\bf Safeguards}
    \item[] Question: Does the paper describe safeguards that have been put in place for responsible release of data or models that have a high risk for misuse (e.g., pre-trained language models, image generators, or scraped datasets)?
    \item[] Answer: \answerNA{}
    \item[] Justification: The dataset contains only indoor tabletop scenes with daily objects. It does not contain faces, personal data, generative models, or any content that poses a misuse risk.

\item {\bf Licenses for existing assets}
    \item[] Question: Are the creators or original owners of assets (e.g., code, data, models), used in the paper, properly credited and are the license and terms of use explicitly mentioned and properly respected?
    \item[] Answer: \answerYes{}
    \item[] Justification: All evaluated models are cited with their original publications and checkpoint sources. The dataset is released under CC BY 4.0; the toolkit under MIT. Overlap with YCB and FewSOL object sets is acknowledged with citations.

\item {\bf New assets}
    \item[] Question: Are new assets introduced in the paper well documented and is the documentation provided alongside the assets?
    \item[] Answer: \answerYes{}
    \item[] Justification: The dataset is accompanied by a Croissant metadata file (machine-readable dataset card), a HuggingFace dataset card, and the \texttt{rpx-benchmark} toolkit with auto-generated API documentation from source-level docstrings. The collection protocol, consent details, privacy considerations, and known biases are documented in the dataset card.

\item {\bf Crowdsourcing and research with human subjects}
    \item[] Question: For crowdsourcing experiments and research with human subjects, does the paper include the full text of instructions given to participants and screenshots, if applicable, as well as details about compensation (if any)?
    \item[] Answer: \answerYes{}
    \item[] Justification: Data capture and annotation were performed by lab personnel (not crowdsourced). The annotation protocol (SAM2 auto-masks with iterative human refinement) is described in Section~\ref{sec:dataset}. All participants provided informed consent for data collection and public release.

\item {\bf Institutional review board (IRB) approvals or equivalent for research with human subjects}
    \item[] Question: Does the paper describe potential risks incurred by study participants, whether such risks were disclosed to the subjects, and whether Institutional Review Board (IRB) approvals (or an equivalent approval/review based on the requirements of your country or institution) were obtained?
    \item[] Answer: \answerNA{}
    \item[] Justification: The study involves only video capture of tabletop objects and human hands during manipulation. No experimental interventions on human subjects were conducted. Participants are lab members who consented to hand appearance in the interaction-phase recordings.

\item {\bf Declaration of LLM usage}
    \item[] Question: Does the paper describe the usage of LLMs if it is an important, original, or non-standard component of the core methods in this research? Note that if the LLM is used only for writing, editing, or formatting purposes and does \emph{not} impact the core methodology, scientific rigor, or originality of the research, declaration is not required.
    \item[] Answer: \answerNA{}
    \item[] Justification: LLMs are not part of the core method. LLMs were used only for writing assistance and code generation during toolkit development, which does not impact the scientific methodology or experimental results.

\end{enumerate}
